\chapter{有理分式}

本章中,我们默认$K$是域.

\section{有理分式的定义}

\begin{definition}{有理分式域}{}
	定义$K((X_{i})_{i\in I})\coloneq\Frac(K[(X_{i})_{i\in I}])$.
	\begin{enumerate}
		\item 我们称$K((X_{i})_{i\in I})$是以$K$中元素为系数,以$X_{i}$为不定元的有理分式域.
		\item 我们称$K((X_{i})_{i\in I})$的元素是以$K$中元素为系数,以$X_{i}$为不定元的有理分式.
		\item 当$I=\left[1,n\right]$时,记$K(X_{1},\ldots,X_{n})\coloneq K((X_{i})_{i\in I})$.
	\end{enumerate}
\end{definition}

\begin{proposition}{整环上的多项式环的分式域}{}
	设$A$是整环,$K=\Frac\left(A\right)$,则有典范$A$-代数同构$A((X_{i})_{i\in I})\simeq K((X_{i})_{i\in I})$.
\end{proposition}

\begin{proposition}{有理分式域的迭代结构}{}
	设$J\subseteq I,K^{\prime}=K\left(\left(X_{i}\right)_{i\in J}\right)$,则有典范$A$-代数同构$K((X_{i})_{i\in I})\simeq K^{\prime}\left(\left(X_{i}\right)_{i\in I\setminus J}\right)$.
\end{proposition}



